\documentclass[aspectratio=169,10pt,t]{beamer}
\usepackage{ac1}
\aulainfo{01}{Apresentação do Curso}{Pensamento lógico e a ideia de algoritmo}

\begin{document}
\capa

\begin{frame}{Roteiro da aula}
  \begin{itemize}
    \item Quem somos e o que esperamos da disciplina
    \item Plano de ensino e avaliação
    \item Por que programar? Por que Python?
    \item Pensamento lógico: o que é um \alert{algoritmo}
    \item Primeiro contato com um programa real
  \end{itemize}
  \vfill
  {\footnotesize Referência: Prefácio do livro \emph{Introdução à Programação e aos Algoritmos} (J.~A.~Ribeiro, LTC, 2019).}
\end{frame}

\begin{frame}{Por que um engenheiro precisa programar?}
  \begin{columns}[T]
    \column{0.55\textwidth}
    \begin{itemize}
      \item Tecnologia está em tudo: carro, geladeira, celular, rede elétrica\ldots
      \item Quem não entende a tecnologia fica apenas como \emph{usuário}
      \item Engenheiros automatizam cálculos, processam dados de medições e simulam sistemas
      \item Programar é uma ferramenta de trabalho, como a calculadora e a planilha
    \end{itemize}
    \column{0.42\textwidth}
    \begin{exampleblock}{Frase do prefácio}
      \itshape ``Vivemos em uma sociedade intensamente dependente da ciência e da tecnologia, em que quase ninguém sabe algo sobre ciência e tecnologia.'' (Carl Sagan)
    \end{exampleblock}
  \end{columns}
\end{frame}

\begin{frame}{Programar se aprende programando}
  \begin{block}{A analogia da culinária}
    Para cozinhar bem é preciso \textbf{teoria} (como os ingredientes se combinam) e \textbf{prática} (ir para a cozinha). Com programação é igual.
  \end{block}
  \begin{itemize}
    \item Ler livros sem programar \(\Rightarrow\) não funciona
    \item Programar sem base \(\Rightarrow\) descobrir sozinho o que outros já sabem há décadas
    \item Vantagem de programar: \alert{feedback imediato}, o erro aparece assim que o programa roda
  \end{itemize}
  \begin{alertblock}{Regra de ouro da disciplina}
    Todo conceito visto em aula deve ser digitado e executado por você.
  \end{alertblock}
\end{frame}

\begin{frame}{Estrutura da disciplina (32+ aulas)}
  \centering
  \small
  \begin{tabular}{@{}p{3.3cm}p{6.7cm}l@{}}
    \toprule
    \textbf{Bloco} & \textbf{Conteúdo} & \textbf{Aulas}\\
    \midrule
    1. Computação essencial & Bases numéricas, arquitetura, algoritmos & 1--4\\
    2. Fundamentos & Tipos, decisão, repetição, funções, módulos, exceções & 5--15\\
    3. Estruturas de dados & Listas, matrizes, dicionários, NumPy, strings, arquivos & 16--25\\
    4. Análise de dados e projeto & Matplotlib, Pandas e projeto final & 26--35\\
    \bottomrule
  \end{tabular}
  \vspace{2ex}

  \footnotesize Provas: \textbf{P1} (aula 15) e \textbf{P2} (aula 25). Projeto final em grupos de 2 ou 3 alunos.\\
  \emph{Engenharia de Computação:} acrescenta os capítulos de Recursão e Ordenação.
\end{frame}

\begin{frame}{Material e ferramentas}
  \begin{itemize}
    \item \textbf{Livro-texto:} Ribeiro, J.~A. \emph{Introdução à Programação e aos Algoritmos}. LTC, 2019.
    \item \textbf{Site do livro:} exercícios com solução comentada e exercícios extras
          \par\href{\siteurl}{\texttt{\siteurl}}
    \item \textbf{Linguagem:} Python 3
    \item \textbf{Ambiente:} Python instalado (Apêndice do livro), editor de texto/IDE e Jupyter para a parte de dados
  \end{itemize}
  \begin{block}{Combinado}
    Faça os exercícios do site \emph{antes} de olhar a solução. Praticar é o que gera autonomia.
  \end{block}
\end{frame}

\begin{frame}{Pensamento lógico}
  Programar é, antes de tudo, \textbf{resolver problemas de forma organizada}.
  \vspace{1ex}
  \begin{enumerate}
    \item Entender o problema (entradas, saídas, restrições)
    \item Dividir em passos pequenos e sem ambiguidade
    \item Só então escrever em uma linguagem de programação
    \item Testar com exemplos e corrigir
  \end{enumerate}
  \vspace{1ex}
  \begin{exampleblock}{Pergunta}
    Como você explicaria a um robô, passo a passo, como atravessar uma rua?
  \end{exampleblock}
\end{frame}

\begin{frame}{O que é um algoritmo?}
  \begin{block}{Definição}
    Sequência \textbf{finita} de passos \textbf{precisos} e \textbf{não ambíguos} que, a partir de entradas, produz uma saída e resolve um problema.
  \end{block}
  \begin{columns}[T]
    \column{0.5\textwidth}
    \textbf{Exemplos do dia a dia}
    \begin{itemize}
      \item Receita de bolo
      \item Trocar um pneu
      \item Achar uma palavra no dicionário
      \item Fazer a divisão ``armada''
    \end{itemize}
    \column{0.5\textwidth}
    \textbf{Propriedades}
    \begin{itemize}
      \item Entrada e saída bem definidas
      \item Cada passo é executável
      \item Termina após um número finito de passos
    \end{itemize}
  \end{columns}
  \vfill
  {\footnotesize O nome vem de Al-Khwarizmi (matemático persa, séc.~IX).}
\end{frame}

\begin{frame}[fragile]{Um algoritmo escrito em português}
  \textbf{Problema:} dado um número, descobrir se ele é par.
  \begin{enumerate}
    \item Leia o número \(n\).
    \item Divida \(n\) por 2 e observe o resto.
    \item Se o resto for 0, diga ``par''.
    \item Caso contrário, diga ``ímpar''.
  \end{enumerate}
  \vspace{1ex}
  \textbf{O mesmo algoritmo em Python:}
\begin{lstlisting}
n = int(input("Digite um número inteiro: "))
if n % 2 == 0:
    print(n, "é par")
else:
    print(n, "é ímpar")
\end{lstlisting}
  \fonte{Exercício 3.4 do livro}
\end{frame}

\begin{frame}[fragile]{Primeiro contato com Python}
  \begin{columns}[T]
    \column{0.52\textwidth}
\begin{lstlisting}
# calcula a área de um triângulo
base = 5
altura = 6
area = (base * altura) / 2
print('A área é', area)
\end{lstlisting}
    \column{0.44\textwidth}
\begin{lstlisting}[style=saida]
A área é 15.0
\end{lstlisting}
    \vspace{1ex}
    \begin{itemize}
      \item \texttt{=} guarda um valor em uma variável
      \item \texttt{print} mostra o resultado
      \item \texttt{\#} inicia um comentário
    \end{itemize}
  \end{columns}
  \vfill\fonte{code/Cap2/prog.2.5\_area\_triangulo\_v3.py}
\end{frame}

\begin{frame}{Atividade em dupla (5 minutos)}
  Escreva em português, com passos numerados, um algoritmo para uma destas tarefas:
  \begin{enumerate}
    \item Fritar um ovo \hfill {\footnotesize (Exercício 2.1)}
    \item Ensinar uma criança de 6 anos a escovar os dentes \hfill {\footnotesize (Exercício 2.2)}
    \item Encontrar o maior entre três números
  \end{enumerate}
  \vspace{1ex}
  Depois troque com outra dupla: o algoritmo deles é \alert{sem ambiguidade}? Dá para executar exatamente como está escrito?
\end{frame}

\exercicios{2}{Exercícios \textbf{2.1} e \textbf{2.2} (algoritmos do cotidiano). Os demais serão vistos na Aula 4.}

\resumo{
  \item Programar = pensar logicamente + praticar muito
  \item Algoritmo: passos finitos, precisos e sem ambiguidade
  \item Usaremos \textbf{Python} e o livro de J.~A.~Ribeiro como base
  \item Próxima aula: como o computador representa informação (binário, hexadecimal, caracteres)
}

\end{document}
